\documentclass[11pt]{article}
\usepackage[margin=1in]{geometry}
\usepackage{amsmath,amssymb,amsthm}

\newtheorem{theorem}{Theorem}
\newtheorem{corollary}{Corollary}
\newtheorem{remark}{Remark}

\title{Intrinsic Coercivity in the ASD Regime}
\author{}
\date{}

\begin{document}
\maketitle

\begin{theorem}[ASD intrinsic spectral gap]
Let $(M^4,g)$ be a compact oriented Riemannian $4$-manifold and let $A$ be an anti-self-dual
Yang--Mills connection on a principal $\mathrm{SU}(2)$-bundle over $M$.
Assume
\[
F_A^+=0
\qquad\text{and}\qquad
W_g^+=0.
\]
If the scalar curvature satisfies
\[
s_g\ge s_{\min}>0,
\]
then the deformation operator $\Lambda_A$ on the Coulomb slice satisfies
\[
\lambda_1(\Lambda_A)\ge \frac{s_{\min}}{4}.
\]
\end{theorem}

\begin{proof}
In the anti-self-dual deformation complex, the Weitzenb\"ock formula on the Coulomb slice has the form
\[
\Lambda_A\alpha=\nabla_A^*\nabla_A\alpha+\frac{s_g}{4}\alpha+W_g^+\!\cdot\alpha+F_A^+\!\cdot\alpha.
\]
Under the hypotheses $F_A^+=0$ and $W_g^+=0$, this reduces to
\[
\Lambda_A\alpha=\nabla_A^*\nabla_A\alpha+\frac{s_g}{4}\alpha.
\]
Taking the $L^2$ inner product with $\alpha$ yields
\[
\langle \Lambda_A\alpha,\alpha\rangle
=
\|\nabla_A\alpha\|_{L^2}^2+\int_M \frac{s_g}{4}|\alpha|^2\,d\mathrm{vol}
\ge
\|\nabla_A\alpha\|_{L^2}^2+\frac{s_{\min}}{4}\|\alpha\|_{L^2}^2.
\]
Hence
\[
\langle \Lambda_A\alpha,\alpha\rangle\ge \frac{s_{\min}}{4}\|\alpha\|_{L^2}^2,
\]
so the Rayleigh quotient gives
\[
\lambda_1(\Lambda_A)\ge \frac{s_{\min}}{4}.
\]
\end{proof}

\begin{corollary}[Round four-sphere]
For the round $S^4$ of radius $1$, one has $s_g=12$ and $W_g^+=0$.
Therefore every ASD connection satisfies
\[
\lambda_1(\Lambda_A)\ge 3.
\]
\end{corollary}

\begin{proof}
Apply the theorem with $s_{\min}=12$.
\end{proof}

\begin{remark}
The statement is intrinsic in the ASD conformally-flat regime.
Without the hypotheses $F_A^+=0$ and $W_g^+=0$, additional curvature terms remain and positivity
requires a separate estimate.
\end{remark}

\end{document}
